% NOTE: The "Default Bibliography Tool" needs to be changed to biber if compiled in TeXstudio

\documentclass[a4paper,9pt]{extarticle}

% Packages
\usepackage[utf8]{inputenc} % For input encoding
\usepackage{geometry} % For page margins
\geometry{a4paper, margin=0.6in} % Set paper size and margins
\usepackage{titlesec} % For section title formatting
\usepackage{enumitem} % For itemized list formatting
\usepackage{hyperref} % For hyperlinks
\usepackage{changepage} %For Indenting

\hypersetup{ % Prettier links
	colorlinks=true,       % True: colored text; False: boxed links
	linkcolor=red,         % Color of internal links (sections, pages, figures)
	citecolor=green,       % Color of library citations
	filecolor=cyan,        % Color of links to local files
	urlcolor=magenta,      % Color of website URLs
	allcolors=blue         % Optional: Sets all link types to one uniform color
}

% Formatting
%\setlist{noitemsep} % Removes item separation
\titleformat{\section}{\Large\bfseries}{\thesection}{1em}{}[\titlerule] % Section title format
\titlespacing*{\section}{0pt}{\baselineskip}{\baselineskip} % Section title spacing
\newenvironment{subs} %Indenting
{\adjustwidth{2em}{0pt}}
{\endadjustwidth}


\titleformat{\subsection}{\large\bfseries}{}{0em}{}
%[\titlerule]
\titlespacing*{\subsection}{0pt}{\baselineskip}{\baselineskip} % Section title spacing

%----------------Bibliography Setup---------
\usepackage[
backend=biber,
style=authoryear,
sorting=ydnt,      % Year (descending), name, title
maxbibnames=4,     % Show first 3 authors
minbibnames=4,
giveninits=true,   % James -> J.
doi=true,
isbn=false,
url=false,
eprint=false
]{biblatex}

\addbibresource{publications.bib}

%----------------Bibliography Customization---------
% Remove quotation marks
\DeclareFieldFormat[article,inproceedings]{title}{#1}

% less whitespace
%\setlength{\bibitemsep}{0.4\baselineskip}
%\setlength{\bibhang}{1.8em}

% italic journal
%DeclareFieldFormat{journaltitle}{\mkbibemph{#1}}

% supress "in:"
%\renewbibmacro{in:}{}

%\ExecuteBibliographyOptions{
%	maxbibnames=4,
%	minbibnames=2
%}

%----------------Bibliography Driver---------
% Declar bibliography environment
\defbibenvironment{bibliography}
{\list
	{}
	{
		\setlength{\leftmargin}{0pt}
		\setlength{\itemindent}{0pt}
		\setlength{\labelwidth}{0pt}
		\setlength{\labelsep}{0pt}
		\setlength{\itemsep}{0.8em}
		\setlength{\parsep}{0pt}
}}
{\endlist}
{\item}

\newlength{\yearwidth}
\setlength{\yearwidth}{3em}

% Article Driver
\DeclareBibliographyDriver{article}{%
	\noindent
	\makebox[\yearwidth][l]{%
		\textbf{\printfield{year}}%
	}%
	\begin{minipage}[t]{\dimexpr\linewidth-\yearwidth\relax}
		\textbf{\printfield{title}}\adddot\space
		
		\printnames{author}\adddot\space
		
		\mkbibemph{\printfield{journaltitle}}%
		\setunit{\space}
		\textbf{\printfield{volume}}
		\iffieldundef{number}
		{}
		{(\printfield{number})}
		%
		\setunit{\addcomma\space}
		\printfield{pages}
		
		\printfield{url}
		\finentry
		
	\end{minipage}
	
	\par\vspace{0.5\baselineskip}
}

% Preprint Driver
\DeclareBibliographyDriver{misc}{%
	\noindent
	\makebox[\yearwidth][l]{%
		\textbf{\printfield{year}}%
	}%
	\begin{minipage}[t]{\dimexpr\linewidth-\yearwidth\relax}
		\textbf{\printfield{title}}\adddot\space
		
		\printnames{author}\adddot\space
		
		\mkbibemph{\printfield{journaltitle}}%
		\setunit{\space}
		\textbf{\printfield{volume}}
		\iffieldundef{number}
		{}
		{(\printfield{number})}
		%
		\setunit{\addcomma\space}
		\printfield{pages}
		\finentry
		
	\end{minipage}
	
	\par\vspace{0.5\baselineskip}
}

%-----------------Command Definitions------------
% Define the CV Entry command
\newcommand{\cventry}[3]{
	\noindent
	\begin{tabular*}{\textwidth}{@{\extracolsep{\fill}}p{2.3cm}p{13cm}}
		\textbf{#1} &
		\textbf{#2}\\
		&
		#3\\[0.8em]
	\end{tabular*}
}

\newcommand{\ccventryy}[4]{
	\noindent
	\begin{tabular*}{\textwidth}{@{\extracolsep{\fill}}p{2.3cm}p{13cm}}
		\textbf{#1} &
		\textbf{#2}\\
		&
		#3\\
		&
		#4\\[0.8em]
	\end{tabular*}
}

\newcommand{\cventryy}[4]{
	\noindent
	\begin{tabular*}{\textwidth}{@{\extracolsep{\fill}}p{0.6cm}p{14.7cm}}
		\textbf{#1} &
		\textbf{#2}\\
		&
		#3\\
		&
		#4\\[0.8em]
	\end{tabular*}
}

\newcommand{\cventryyy}[5]{
	\noindent
	\begin{tabular*}{\textwidth}{@{\extracolsep{\fill}}p{0.6cm}p{14.7cm}}
		\textbf{#1} &
		\textbf{#2}\\
		&
		#3\\
		&
		#4\\
		&
		#5\\[0.8em]
	\end{tabular*}
}

% Define the CV Item command, displays a bullet point
\newcommand{\cvitem}[1]{
	\hspace{2.3cm}$\bullet$ #1\\
}

%-------------------------------------------------------------------------------
% Begin document
%-------------------------------------------------------------------------------
\begin{document}
	
	% Disable page numbers
	\pagestyle{empty}
	
	% Header
	\begin{center}
		{\LARGE\bfseries James R. Bounds, PhD}
		
		\vspace{3pt}
		
		Curriculum Vitae
		
		\vspace{4pt}
		
		Austin, Texas \quad$\cdot$\quad
		\href{mailto:jamesbounds07@gmail.com}{jamesbounds07@gmail.com} % | \href{https://www.linkedin.com/in/loremipsum}{Website or LinkedIn} % Contact info
	\end{center}
	
	\section*{EXECUTIVE SUMMARY}
	Experimental Physicist (PhD) and Hardware/Software Engineer with 10+ years of experience designing precision optical instruments, custom analog electronics, and scientific automation software. Expert in laser spectroscopy, frequency combs, and high-performance numerical computing.
	
	%------------------------
	\section*{EDUCATION \& EMPLOYMENT HISTORY}
	\ccventryy%y
	{2015}
	{B.S. Physics}
	{Texas A\&M University}
	{Department of Physics and Astronomy}
%	{GPA: 3.19}
	
	\ccventryy
	{2012--2015}
	{Undergraduate Research Assistant}
	{Texas A\&M University}
	{Advisor: Prof. Hans Schuessler}
	
	\ccventryy
	{2015-2016}
	{IT Technician}
	{Systek Computing}
	{College Station, Texas}
	
	\ccventryy
	{2016--2025}
	{Graduate Research Assistant}
	{Texas A\&M University}
	{Advisor: Prof. Hans Schuessler}
	
	\ccventryy
	{2016-2025}
	{Visiting Scientist (Summer terms)}
	{Max Planck Institute of Quantum Optics}
	{Supervisor: Prof. Thomas Udem}
	
	\ccventryy
	{2025--2026}
	{Graduate Teaching Assistant}
	{Texas A\&M University}
	{Advisor: Prof. Hans Schuessler, Prof. Alexei Sokolov}
	
	\ccventryy
	{2026}
	{PhD. Physics}
	{Texas A\&M University}
	{Advisor: Prof. Alexei Sokolov}
	
	\section*{RESEARCH EXPERIENCE}
	\begin{itemize}
		\item \textbf{Designed, assembled, and aligned high-precision optical setups:} \newline
		Including Cavity Ring-down Spectroscopy (CRDS), Vernier spectroscopy, and Doppler-free saturated absorption systems. These systems achieved sub-ppb sensitivity in trace-gas detection.
		
		\item \textbf{Theoretical modeling and optimization of complex optical systems:} \newline
		Built Gaussian beam simulation framework and custom ray/wave propagation tools in Python and C to validate beam geometry and cavity coupling efficiency.
		
		\item \textbf{Engineered custom optoelectronic and RF detection hardware:} \newline
		Full life-cycle development of custom high-bandwidth amplified photodetectors, analog signal conditioning circuits, and mixed digital/analog devices for the optimization of experimental signals. This development spans the conception, schematic, KiCAD PCB layout, assembly, and enclosure design.
		
		\item \textbf{Calibrated and characterized mid-IR and near-IR laser systems:} \newline
		Used optical frequency combs, custom-built wavemeters and interferometers for ultra-high-resolution wavelength metrology, optical frequency metrology, and hyperfine structure determinations.
		
		\item \textbf{Developed full-stack laboratory automation and data acquisition (DAQ) software:} \newline
		Development spans from hardware-level FPGA fabric design for accelerated real-time processing, OS-level drivers written in C, to the high level interface written in Python and QT.
		
		\item \textbf{Constructed precision optomechanical assemblies and vacuum interface systems:} \newline
		utilizing CAD software (Autodesk Inventor), 3D printing, and hands-on lathe/milling fabrication to ensure structural stability and low thermal drift of optical and vacuum-compatible mounts.
		
		\item \textbf{Executed optical fiber characterization and assembly:}
		\newline
		including single-mode and polarization-maintaining fiber splicing, free-space-to-fiber coupling, and loss minimization across visible and mid-IR wavelengths.
		
		\item \textbf{Led end-to-end experimental error analysis and validation:} \newline
		Processing raw spectral measurements against line-list databases (eg. HITRAN) using custom non-linear fitting and numerical optimization algorithms, in addition to error quantification.
		
		\item \textbf{Collaborated across international research divisions:} \newline
		(including Max Planck Institute of Quantum Optics), delivering hardware prototypes and presentation deliverables on tight timelines.
		
	\end{itemize}
	
%	\section*{RESEARCH EXPERIENCE}
%	\begin{itemize}
%		\item Design and construction of high precision scientific instruments
%		\item Conception of novel data analysis algorithms
%		\item Implementation of laser spectroscopy and cavity-enhanced absorption techniques
%		\item Characterization of breath samples using laser spectroscopy
%		\item High-resolution measurement of molecular absorption cross section
%		\item Retrieval of hyperfine constants from doppler-free spectra
%		\item Demonstration of novel wavelength and optical frequency calibration techniques
%		\item Optical system design and alignment (CRDS, Vernier spectroscopy, interferometry)
%		\item Scientific instrument design (high precision spectrometers and wavemeters, novel measurement platforms)
%		\item Scientific software development (full-stack, numerical optimization, instrument automation, GUI design)
%		\item Electronics and instrumentation (DAQ systems, photodetectors, control electronics)
%		\item Experimental characterization, uncertainty analysis, and validation against molecular databases
%	\end{itemize}

	\section*{PROJECTS AND HARDWARE SYSTEMS}
	\begin{itemize}
		\item \textbf{Optical Breath Analysis Platforms:} Complete design, fabrication, assembly and alignment of both cavity-ringdown spectrometers and Vernier spectrometers. Both systems are easily capable of sub-1ppb trace gas detection in human breath samples.
		
		\item \textbf{Wavelength Metrology Platform:} Construction of a real-time wavelength analysis instrument capable of precision of less than 1 part in $10^8$, aided by two complementary algorithms invented in my dissertation for the analysis of instrument data.
		
		\item \textbf{Ultra-high precision IR spectrometer:} Conception, design, and implementation of a saturated absorption spectrometer with on-board optical frequency calibration using a comb. This was used to measure the hyperfine spectrum of metastable xenon with accuracy competitive with much more complex systems.
		
	\end{itemize}
	
	%------------------------
	% Publications Section
	\section*{PUBLICATIONS}
	
	\nocite{*}
	
	
	\subsection*{Peer Reviewed}
	
	\printbibliography[heading=none,keyword=journal]
	
	\subsection*{Preprint}
	\printbibliography[heading=none,keyword=preprint]
	
	\section*{CONFERENCE PROCEEDINGS}
	
	\cventryy
	{2023}
	{Hyperfine Structure of Xe by Frequency Comb Calibrated Saturation and Collinear Fast Beam Spectroscopies}
	{31st Annual International workshop, Laser Physics 2023}
	{H.Schuessler, J. Bounds, N. Schnoor, A. Kolomenskii, R. Trainham, M. Manard.}
	
	\cventryy
	{2023}
	{Local and Remote Measurements of Atmospheric Trace Gases}
	{31st Annual International workshop, Laser Physics 2023}
	{A. Kolomenskii, C. Rodriguez, J. Bounds, E. Aluauee, H. Schuessler.}
	
	\cventryyy
	{2021}
	{Broadband Frequency Comb Vernier Spectroscopy In The Mid-IR With Video-recording Retrieval Of Absorption Spectrum}
	{International Symposium on Molecular Spectroscopy}
	{H. Schuessler, J. Bounds, A. Kolomenskii}
	{\href{http://hdl.handle.net/2142/111474}{DOI: 10.15278/isms.2021.FF12}}
	
	\cventryyy
	{2021}
	{Spectroscopic Detection Of Methane at a ppt Sensitivity Level in Mid-IR with a Long-path Multipass Cell}
	{International Symposium on Molecular Spectroscopy}
	{J. Xia, C. Rodriguez, J. Bounds, A. Kolomenskii, H. Schuessler}
	{\href{http://hdl.handle.net/2142/111469}{DOI: 10.15278/isms.2021.RF10}}
	
	\cventryyy
	{2018}
	{Frequency comb vernier spectroscopy of methane in the mid-IR with temporal retrieval of comb lines}
	{International Symposium on Molecular Spectroscopy}
	{James R Bounds, F. Zhu, A. Kolomenskii, H. Schuessler.}
	{\href{http://hdl.handle.net/2142/100480}{DOI: 10.15278/isms.2018.WA05}}
	
	\cventryyy
	{2018}
	{Acetone and methane detection with wavelength modulation spectroscopy in the near- and mid-IR}
	{International Symposium on Molecular Spectroscopy}
	{J. Xia, F. Zhu, J. Bounds, A. Kolomenskii, H. Schuessler.}
	{\href{http://hdl.handle.net/2142/100674}{DOI: 10.15278/isms.2018.RH03}}
	
	\section*{POSTER PRESENTATIONS}
	\cventryy
	{2025}
	{Hyperfine structure and isotope shifts of xenon measured for near-infrared transitions with Doppler-free saturated absorption spectroscopy}
	{SIBOR Workshop}
	{Cook's Branch Conservancy, Montgomery, Texas}
	
	\section*{ORAL PRESENTATIONS}
	
	\cventryy
	{2024}
	{Frequency Comb Calibration of the Hyperfine Structure of Metastable Xenon}
	{SIBOR Workshop}
	{Cook's Branch Conservancy, Montgomery, Texas}
	
	\cventryy
	{2023}
	{Frequency Comb Calibration of the Hyperfine Structure of Metastable Xenon}
	{SIBOR Workshop}
	{Cook's Branch Conservancy, Montgomery, Texas}
	
	\cventryy
	{2022}
	{Metastable Xenon Hyperfine Structure Measurement}
	{SIBOR Workshop}
	{Cook's Branch Conservancy, Montgomery, Texas}
	
	\cventryy
	{2019}
	{Vernier Spectroscopy and Spectroscopic Detection of Biomarkers in Exhaled Breath}
	{Workshop of Max Planck Institute of Quantum Optics Laser Spectroscopy Division}
	{Schloss Ringberg, Kreuth, Germany}
	
	\section*{OPEN SOURCE PROJECTS}
	
	Github: \url{https://github.com/jaymz07}
	
	\begin{itemize}
		\item \textbf{Battery-Powered, Amplified Photodetector} \newline
			\null\qquad \url{https://github.com/jaymz07/Photodetector-Battery}
			
		\item \textbf{Quantum Tunneling Demo Written in Java} \newline
		\null\qquad \url{https://github.com/jaymz07/quantum-tunneling}
		
		\item \textbf{Gaussian Beam Simulator GUI} \newline
		\null\qquad \url{https://github.com/jaymz07/Gaussian-Beam-Sim-Qt}
		
		\item \textbf{Grotrian Diagram Generator} \newline
		\null\qquad \url{https://github.com/jaymz07/Grotrian}
	\end{itemize}
	
	\section*{TECHNICAL SKILLS}
	
	\begin{description}[font=\bfseries, leftmargin=0pt, labelsep=0.5em, itemsep=0.4em]
		\item[Optics \& Photonics:] Precision Free-Space Alignment, Optical Frequency Combs, Cavity-Enhanced Spectroscopy (CRDS, Vernier), Dual-Comb Spectroscopy, Optical Fiber Splicing \& Coupling, Interferometry, Wavelength Metrology, Multi-Pass Cells, Vacuum Optical Interfaces.
		
		\item[Hardware, EDA \& Electronics:] Analog/Digital/Mixed-Circuit Design, Schematic Capture \& PCB Layout (KiCAD), Circuit Simulation (NI Multisim, QUCS), FPGA Fabric Design (Verilog), Microcontrollers \& Embedded Systems, Photodetector Optimization, Signal Conditioning, PID Controllers, VNAs, Oscilloscopes.
		
		\item[Instrumentation \& DAQ:] Hardware Protocol Automation (GPIB, RS232, USB, Ethernet, SPI, $I^2C$, PyVISA), Custom DAQ Systems, FPGA Acceleration, GUI Development (Qt/PyQt, Tkinter, LabVIEW).
		
		\item[Programming Languages:] Python, C, C++, Cython, Verilog, MATLAB, Java, Bash, Shell Scripting, Git, \LaTeX.
		
		\item[Scientific Computing \& Modeling:] NumPy, SciPy, PyTorch, Numba, OpenCV, Scikit-learn, Pandas, Slurm HPC Cluster Management, Ray-Tracing \& Gaussian Beam Modeling.
		
		\item[Mechanical:] 3D and CAD Modeling (Autodesk Inventor, Blender, OpenGL, PyVista), Precision Machining (Vertical Mill, Precision Lathe), 3D Printing / Additive Manufacturing, Ultra-High Vacuum (UHV) Systems.
		
		\item[Systems \& Infrastructure:] Linux Systems Administration (openSUSE, Debian, Arch, Fedora, RHEL/CentOS), SSH Server Configuration, Network Infrastructure, X11 Forwarding.
	\end{description}
	
\end{document}
%-------------------------------------------------------------------------------
%-------------------------------------------------------------------------------

